\documentclass{article}
\title{확률 컴퓨터란 무엇일까?}
\usepackage{kotex}
\author{윤도현}
\begin{document}
	\maketitle
	\section{개요}
	아주 멋지게도 인류는 전기를 다양한 소자를 이용하여 스위칭하는 법을 알아낸 이후 이를 응용하여 온갖 다양하고 개성 넘치는 물건을 만들어냈다. 이 중 가장 멋지다고 평가할 만한 물건인 컴퓨터는 전 인류 역사를 통틀어 가장 혁신적이라 평가할 만한 물건이다. 이미 컴퓨터는 기상, 금융, 과학, 의료 등 이젠 쓰이지 않는 분야를 찾기 더 힘들어질 정도로 우리의 삶 깊숙이 자리 잡았으며 더 이상 때려야 땔 수 없는 관계를 구축했다. 하지만 이에 따라 발생하는 엄청난 양의 에너지 사용과 기타 등등의 문제도 함께 동반되어 나타났고 이를 해결하기 위해 TSMC, 삼성 등의 기업은 공정 미세화를 지속적으로 연구 및 개발하였으나 물리적인 한계에 직면하였다. 따라서 우린 새로운 방식의 고효율 컴퓨팅법을 통하여 타게 하고자 하며 그 방법 중 하나인 확률컴퓨팅을 소개하고자 한다.
	\section{작동 원리}
	우선 확률 컴퓨터의 가장 기초적인 연산 단위인 확률 비트(P-bit)에 대하여 알아보자. 
	\subsection{P-bit란 무엇일까?}
	P-bit란 간단히 말해 조절 가능한 확률을 가진 무작위의 값을 가질 수 있는 비트이다. 이게 뭔 미친 소리인가 싶을 수 있다. 그게 매우 정상적인 반응이다. 우린 이것을 이해하기 위해 아주 간단한 예시를 통해 P-bit에 대한 개념을 살펴볼 것이다.
	\subsubsection{동전 던지기}
	동전 던지기를 할 때 가장 이상적인 상황에서 앞면과 뒷면이 나올 확률은 각각 50:50이다. 이러한 조건에 계속해서 동전을 던진다면 우린 결국 50:50에 가까운 값으로 앞면과 뒷면이 나온다는 사실을 알 수 있을 것이다. 하지만 한쪽 면에 무게를 추가한다거나 뺀다든지 하는 편향을 주게 된다면 어떻게 될까? 이런 식으로 편향을 준다면 해당 면은 덜 나오거나 더 많이 나올 수 있을 것이다. P-bit는 이것과 동일한 개념으로 동작한다. P-bit는 아무런 입력이 없는 상황에서는 일정한 주기로 0과 1중 하나의 값을 50:50의 확률로 가진다. 하지만 입력값을 주어 둘 중 하나에 편향을 주게 된다면 동전 던지기와 마찬가지로 특정한 값이 더 많이 나오거나 더 적게 나오게 될 것이다. 이것이 P-bit가 동작하는 방식이다.
	\subsection{그럼 이걸로 어떻게 계산할까?}
	이것만 본다면 "대체 이걸로 어떻게 수학 연산을 한다는 거지?"란 생각이 들 수 있다. 매우 당연한 반응이다. 이걸 이해하기 위해 우린 간단한 문제 예시를 들어 살펴볼 것이다.
	\subsubsection{문제}
	우린 P-bit를 이용하여 
	$$2X = 6$$
	요런 문제를 풀어보며 이 괴상한 컴퓨터를 이해하려 시도할 것이다.	
	\subsubsection{문제 변환하기}
	P-bit를 사용하는 확률 컴퓨터를 이용하여 풀려면 우선 문제를 확률 컴퓨터에서 사용하기 편한 형태로 변환해야 한다. 우선 확률 컴퓨터에서 원활히 사용할 수 있는 형태로 변환해 보자
	\newline
	\newline
	식
	$$
	2X = 6
	$$
	를 에너지로 변환하여
	$$
	E = (2X-6)^2
	$$
	형태로 변환할 수 있다. 이후 이것을 확률 컴퓨터에 넣게 된다. 물론 이것은 매우 단순화된 과정이고 여러 과정이 더 있지만, 이 문서의 목표가 '대략 이렇게 동작한다'를 설명하는 것이기에 생략한다.
	\subsubsection{그럼 이제 어떻게 되나?}
	이런 식으로 변환된 문제를 받은 확률 컴퓨터는, 내부에 존재하는 P-bit들로 이루어진 네트워크를 이용하여 문제를 해결한다. 처음에는 P-bit 네트워크에서는 이러한 문제를 받으면 모든 값을 떠다니며 아무 숫자나 내뱉지만 문제가 올바르게 정의되었다는 전제하에 경우가 늘어날수록 특정한 값이 머무르는 경우가 늘어나게 된다. 이는 P-bit 네트워크가 상호작용을 하며 발생한 것으로 우리는 이를 통하여 연산할 수 있으며 해당 과정은 인간의 개입이 불필요하다. 이러한 과정은 양자컴퓨터와 매우 유사한 방식으로 동작하며 양자컴퓨터 대비 에너지 소모량 등의 부분에서 우위를 점하지 못하지만, 온도 유지, 초전도체 등의 불필요하다는 장점이 존재한다.
	\section{그럼 이걸로 모든 문제를 해결할 수 있나?}
	사실 모든 문제에 이것을 사용할 수는 없다. 여러 가지 종류의 여러 문제 중 이러한 방식으로 동작하는 컴퓨터가 잘 풀 수 있는 문제가 있고 없는 문제가 있다. 따라서 모든 문제를 빠르게 해결하는 게 아닌 특정한 문제를 빠르게 해결하는 용도로 쓰일 수 있다.
	\section{끝맺음}
	많은 사람이 양자컴퓨터에 대하여 들어는 봤겠지만, 확률 컴퓨터란 기계에 대해서는 들어보지 못하였을 것이다. 하지만 양자컴퓨터보다 성능은 떨어지나, 구현이 상대적으로 간단한 이 기계의 경우 양자컴퓨터보다 이른 시일 내로 눈에 띄는 결과를 만들어낼 수 있을 것으로 생각한다. 따라서 양자컴퓨터뿐만 아닌 확률 컴퓨터라는 기계도 있다는 사실을 알아주길 바라며 더 자세한 내용은 기타 자료를 참고 바란다.
\end{document}